\subsection{Core Functional Analytic Theorems}

% **Role:** infinite-dimensional linear algebra powered by analysis.
%
% Hahn--Banach, uniform boundedness, open mapping, closed graph theorem.

The core theorems explain why complete normed spaces behave rigidly under linear maps. Baire category turns pointwise control, surjectivity, and closed graphs into operator statements.

\begin{theorem} \label{thm:hahn-banach-extension} \hlt{Hahn--Banach Extension Theorem}\\
Let $\mathcal{X}$ be a real normed space, let $M\subseteq\mathcal{X}$ be a linear subspace, and let $p:\mathcal{X}\to\R$ be sublinear. If $\varphi:M\to\R$ is linear and $\varphi(x)\leq p(x)$ for all $x\in M$, then $\varphi$ extends to a linear functional $\Phi:\mathcal{X}\to\R$ with $\Phi(x)\leq p(x)$ for all $x\in\mathcal{X}$. In the normed complex form, every bounded functional on a subspace extends with the same norm.
\end{theorem}

\begin{theorem} \label{thm:positive-riesz-representation-c-interval} \hlt{Positive Riesz Representation on $C([a,b])$}\\
Let $J=[a,b]$ and let $G\in C(J)^*$ be positive, i.e. $f\geq0$ on $J$ implies $G(f)\geq0$.\\
Then $\exists$ a finite positive Borel measure $\gamma$ on $\R$ with $\operatorname{supp}\gamma\subseteq J$ s.t. $G(f)=\int_J f\,d\gamma$ $\forall f\in C(J)$.\\
Moreover, $\normop{G}=G(\mathbf{1}_J)=\gamma(J)$.
\begin{proof}
For $a\leq t<b$ and $n$ large, set
\begin{align}
\varphi_{t,n}(s)\defeq
\begin{cases}
1, & a\leq s\leq t,\\
n(t+1/n-s), & t<s<t+1/n,\\
0, & t+1/n\leq s\leq b.
\end{cases} \nonumber
\end{align}
Then $\varphi_{t,n}\downarrow \mathbf{1}_{[a,t]}$ pointwise, hence $G(\varphi_{t,n})$ decreases.\\
Define $g(t)\defeq\lim_{n\to\infty}G(\varphi_{t,n})$; $g(t)=0$ for $t<a$; and $g(t)=G(\mathbf{1}_J)$ for $t \geq b$.\\
Ramp comparisons give $g$ increasing and right-continuous.\\
By Theorem \ref{thm:borel-stieltjes-measure-generation}, $g$ induces a finite Borel--Stieltjes measure $\gamma$ with $\operatorname{supp}\gamma\subseteq J$ and $\gamma(J)=g(b)=G(\mathbf{1}_J)$.\\
Fix $f\in C(J)$ and $\varepsilon>0$. Choose $P:a=t_0<\cdots<t_m=b$ with oscillation $\omega_f(P)\leq\varepsilon$.\\
Put $\Delta g_k\defeq g(t_k)-g(t_{k-1})$. For $n$ large, ramp increments $\varphi_{t_k,n}-\varphi_{t_{k-1},n}$ give continuous approximant $f_{P,n}$ whose $G$-increments approximate $\sum_k f(t_k)\Delta g_k$ within $\varepsilon$, while step approximant $s_P$ satisfies $\norm{f-s_P}_{L^1(\gamma)}\leq\varepsilon\gamma(J)$.\\
Hence $\left|G(f)-\int_Jf\,d\gamma\right|\leq \varepsilon(2\normop{G}+1)$ after absorbing constants into $\varepsilon$. Letting $\varepsilon\downarrow0$ gives $G(f)=\int_Jf\,d\gamma$.\\
Finally, positivity gives $\abs{G(f)}\leq G(\mathbf{1}_J)\norm{f}_\infty=\gamma(J)\norm{f}_\infty$, and $f=\mathbf{1}_J$ gives $\normop{G}\geq\gamma(J)$.\\
\hlt{Structural Commentary}: Positivity turns ramp limits into interval masses, hence a Stieltjes measure.\\
\hlt{Key Insight}: Ramps turn order on $C(J)$ into interval masses.
\end{proof}
\end{theorem}

\begin{corollary} \label{cor:riesz-representation-by-charges-conditional} \hlt{Signed Riesz Representation by a Charge}\\
Let $J=[a,b]$ and let $G\in C(J)^*$ admit $G=G^+-G^-$, $G^\pm\in C(J)^*$; for $G^\pm\geq0$.\\
Then $\exists$ a charge $\lambda:\mathcal{B}(\R)\to\R$ with $\operatorname{supp}\lambda\subseteq J$ such that $G(f)=\int_J f\,d\lambda$ $\forall f\in C(J)$.
\begin{proof} By Theorem \ref{thm:positive-riesz-representation-c-interval}, $\exists$ finite positive Borel measures $\gamma^\pm$ with $\operatorname{supp}\gamma^\pm\subseteq J$; $G^\pm(f)=\int_J f\,d\gamma^\pm$ $\forall f\in C(J)$.\\
Set $\lambda\defeq\gamma^+-\gamma^-$, i.e. $\lambda(E)\defeq\gamma^+(E)-\gamma^-(E)$ for $E\in\mathcal{B}(\R)$. Then $\lambda$ is a charge by Proposition \ref{prop:charges-linear-combinations-and-variation}, $\operatorname{supp}\lambda\subseteq J$, and $G(f)=G^+(f)-G^-(f)=\int_J f\,d\gamma^+-\int_J f\,d\gamma^-=\int_J f\,d\lambda$.\\
\hlt{Structural Commentary}: Positive representations subtract to a charge representation.\\
\hlt{Key Insight}: Signed representation is positive representation plus subtraction.
\end{proof}
\end{corollary}

\begin{theorem} \label{thm:uniform-boundedness-principle} \hlt{Uniform Boundedness Principle}\\
Let $\mathcal{X}$ be Banach and let $\mathcal{Y}$ be normed. If $\mathcal{F}\subseteq\LL(\mathcal{X},\mathcal{Y})$ is pointwise bounded, meaning $\sup_{T\in\mathcal{F}}\norm{Tx}<\infty$ for every $x\in\mathcal{X}$, then $\sup_{T\in\mathcal{F}}\normop{T}<\infty$.
\end{theorem}

\begin{theorem} \label{thm:open-mapping-theorem} \hlt{Open Mapping Theorem}\\
A surjective bounded linear map $T:\mathcal{X}\to\mathcal{Y}$ between Banach spaces is open. Equivalently, $T$ maps neighborhoods of $0$ onto neighborhoods of $0$.
\end{theorem}

\begin{theorem} \label{thm:closed-graph-theorem} \hlt{Closed Graph Theorem}\\
Let $\mathcal{X}$ and $\mathcal{Y}$ be Banach spaces. A linear map $T:\mathcal{X}\to\mathcal{Y}$ is bounded if and only if its graph
\[
\Gamma(T)\defeq\{(x,Tx):x\in\mathcal{X}\}
\]
is closed in $\mathcal{X}\times\mathcal{Y}$.
\end{theorem}
